\section{施瓦兹函数类，时限与带限}\label{sec:Schwartz_Functions}

\subsection{施瓦兹函数在傅里叶变换下的封闭性}\label{sub:fourier_close}
在\ref{sub:schwartz}中，我们介绍了施瓦兹函数类及其对偶空间分布，并直接假定了施瓦兹函数类在傅里叶变换下的封闭性。这一小节将填补这个逻辑漏洞，并介绍施瓦兹函数类的一些其他好处。由于不涉及具体的计算，简洁起见，我们将在本节使用角频率$\omega$。

首先，可以根据施瓦兹函数类的定义\[\mathcal{S} =\left\{\varphi\in C^{\infty}(\mathbb{R} )\left|\lim_{|x|\to\infty}|x|^m\varphi^{(n)}(x)=0,\forall m,n\in\mathbb{N} \right.\right\}\]给出它们在$|x|\to\infty$时的估计，从而施瓦兹函数$\varphi$必满足$\int_{-\infty}^{\infty}|\varphi^{(n)}(x)|\,dx<\infty,n\in\mathbb{N}$，即施瓦兹函数类是$L^1(\mathbb{R} )$的子集，于是\begin{align*}
    \int_{-\infty}^{\infty}\varphi(t)\mathrm{e}^{-\mathrm{i}\omega t}\,dt&\leq \int_{-\infty}^{\infty}|\varphi(t)||\mathrm{e}^{-\mathrm{i}\omega t}|\,dt\\
    &=\int_{-\infty}^{\infty}|\varphi(t)|\,dt<\infty
\end{align*}
即施瓦兹函数的傅里叶变换是良定义的，傅里叶逆变换同理。

显然，施瓦兹函数在求导运算下是封闭的，即$\varphi\in\mathcal{S}\implies \varphi'\in\mathcal{S}$，并且乘以$x$的运算也是封闭的，即$\varphi\in\mathcal{S}\implies x\varphi\in\mathcal{S}$。利用这两个性质，可以证明施瓦兹函数类在傅里叶变换下的封闭性：
\begin{theorem}[施瓦兹函数类在傅里叶变换下的封闭性]
    设$\varphi\in\mathcal{S}$，则$\mathcal{F} \varphi,\mathcal{F} ^{-1} \varphi\in\mathcal{S}$。
\end{theorem}
\begin{proof}
    只需证明$\lim_{|\omega|\to\infty}|\omega|^m \mathcal{F} \varphi^{(n)}(\omega)=0,\forall m,n\in\mathbb{N}$。由傅里叶变换的微分性质（见\ref{sub:general_distributions}），
    \begin{align*}
    \mathcal{F} [\varphi'](\omega)&=\mathrm{i}\omega \mathcal{F} \varphi(\omega)\\
    \mathcal{F} [\varphi^{(n)}](\omega)&=(\mathrm{i}\omega)^n \mathcal{F} \varphi(\omega)
\end{align*}
根据黎曼-勒贝格引理（见\ref{sub:riemman_lebesgue_lemma}），\[\varphi^{(n)}\in L^1(\mathbb{R})\implies\lim_{|\omega|\to\infty}\mathcal{F}\left[ \varphi^{(n)}(\omega)\right]=0\]
即\[\lim_{|\omega|\to\infty}(\mathrm{i}\omega)^n \mathcal{F} \varphi(\omega)\to 0\]
故$\mathcal{F} \varphi(\omega)\in\mathcal{S} $，同理可证明$\mathcal{F}^{-1} \varphi(\omega)\in\mathcal{S} $。
\end{proof}

\subsection{时限与带限的关系}\label{sub:time_freq_limit}
我们来验证，紧支函数的傅里叶变换不再是紧支函数，从而一开始介绍的$\mathcal{C} $在傅里叶变换下不封闭。事实上，非零函数不可能既时限，又带限，即：\begin{theorem}
设$f\in L^1(\mathbb{R}),f\neq 0$，如果$f$在时域上紧支，则其傅里叶变换$\mathcal{F} f(\omega)$不可能在频域上紧支，反之亦然。
    
\end{theorem}

\begin{proof}
我们假设$\text{supp}\ f\in[-T/2,T/2],\text{supp}\ \mathcal{F} f\in[-\Omega/2,\Omega/2]$，证明$f=0$。

\textbf{法一}：
\begin{align*}
    f(t)&=\frac{1}{2\pi}\int_{-\infty}^{\infty}\mathcal{F} f(\omega)\mathrm{e}^{\mathrm{i}\omega t}\,d\omega\\
    &=\frac{1}{2\pi}\int_{-\Omega/2}^{\Omega/2}\mathcal{F} f(\omega)\mathrm{e}^{\mathrm{i}\omega t}\,d\omega
\end{align*}
由于$\text{supp}\ f\in[-T/2,T/2]$，所以$f(t)$在$|t|>T/2$时恒为0，故\[\int_{-\Omega/2}^{\Omega/2}\mathcal{F} f(\omega)\mathrm{e}^{\mathrm{i}\omega t}\,d\omega=0,\forall |t|>T/2\]
对以上恒等式两边同时求n阶导数，得到\[\int_{-\Omega/2}^{\Omega/2}(\mathrm{i}\omega)^n \mathcal{F} f(\omega)\mathrm{e}^{\mathrm{i}\omega t}\,d\omega=0,\forall |t|>T/2\]
取定$t_0>T/2$，有\begin{align*}
    f(t)=&\frac{1}{2\pi}\int_{-\Omega/2}^{\Omega/2}\mathcal{F} f(\omega)\mathrm{e}^{\mathrm{i}\omega t_0}\mathrm{e}^{\mathrm{i}\omega (t-t_0)}\,d\omega
\end{align*}
将$\mathrm{e}^{\mathrm{i}\omega (t-t_0)}$展开成幂级数：\begin{align*}
    \mathrm{e}^{\mathrm{i}\omega (t-t_0)}=&\sum_{n=0}^{\infty}\frac{(\mathrm{i}\omega (t-t_0))^n}{n!}
\end{align*}
则该幂级数在$\mathbb{R}$上收敛且内闭一致收敛，特别地，在有限区间$[-\Omega/2,\Omega/2]$上一致收敛，因此可以将求和与积分换序，得到\begin{align*}
    f(t)=&\frac{1}{2\pi}\int_{-\Omega/2}^{\Omega/2}\mathcal{F} f(\omega)\mathrm{e}^{\mathrm{i}\omega t_0}\sum_{n=0}^{\infty}\frac{(\mathrm{i}\omega (t-t_0))^n}{n!}\,d\omega\\
    =&\frac{1}{2\pi}\sum_{n=0}^{\infty}\frac{((t-t_0))^n}{n!}\int_{-\Omega/2}^{\Omega/2}(\mathrm{i}\omega)^n \mathcal{F} f(\omega)\mathrm{e}^{\mathrm{i}\omega t_0}\,d\omega=0
\end{align*}
这表明只有恒零函数才能做到时域、频域均紧支。

\textbf{法二}：利用解析延拓的唯一性。设$\text{supp}\ f\in[-T/2,T/2]$，则\begin{align*}
    \mathcal{F} f(\omega)&=\int_{-\infty}^{\infty}f(t)\mathrm{e}^{-\mathrm{i}\omega t}\,dt\\
    &=\int_{-T/2}^{T/2}f(t)\mathrm{e}^{-\mathrm{i}\omega t}\,dt
\end{align*}
这是一个可导的实变函数，可以将其做全纯延拓。我们在\ref{sub:introduction_of_laplace}中已经介绍过拉普拉斯变换，不过那时使用的是单边拉普拉斯变换，这里为了涵盖$\text{supp}\ f\in[-T/2,T/2]$的情形，我们使用双边拉普拉斯变换：
\begin{align*}
    F(s)=\int_{-\infty}^{\infty}f(t)\mathrm{e}^{-st}\,dt=\int_{-T/2}^{T/2}f(t)\mathrm{e}^{-st}\,dt,F(\mathrm{i}\omega)=\mathcal{F} f(\omega)
\end{align*}
它同样具备全纯性。由于积分区间有限，$f\in\mathcal{C} $光滑，以上积分总是存在的，并且可以对$s$求导，从而也是全纯的：\[F'(s)=-\int_{-T/2}^{T/2}tf(t)\mathrm{e}^{-st}\,dt\]
即$F(s)$是整函数。但是，$F(s)=0,\forall s=\mathrm{i}\omega,|\omega|> \Omega/2$，由唯一延拓定理，$F(s)=0,\forall s\in\mathbb{C}$，因此$f(t)=0$。
\end{proof}

%时域频域“带宽”积，海森堡不确定性原理
必须指出，实际应用中我们总是考虑时域、频域都近似紧支的函数，指出时限、带限不共存，只是理论需要。事实上，时域上的“宽度”与频域上的“宽度”之积存在一个下界，这一点从\ref{sub:fourier_operation_properties}中提出的伸缩定理就可以看出：
\lr{
    \mathcal{F} [f(at)](\xi)=\frac{1}{|a|}\mathcal{F} f\left(\frac{\xi}{a}\right)
}{
    \mathcal{F} [f(at)](\omega)=\frac{1}{|a|}\mathcal{F} f\left(\frac{\omega}{a}\right)
}
随着$a$的增大，时域函数$f(at)$变得越来越集中，而频域函数$\mathcal{F} f\left(\frac{x}{a}\right)$则变得越来越分散，反之亦然。
我们可以用方差来定量描述这一关系。
\begin{theorem}
    时域信号、频域信号的方差之积满足不等式
    \[\sigma(f)\sigma(F)\geq \frac{1}{4\pi}\]
    其中$f\in L^1(\mathbb{R} )\cap L^2(\mathbb{R} )$，$F(\xi)=\int_{-\infty}^{\infty}f(t)\mathrm{e}^{-2\pi \mathrm{i}\xi t}\,dt$，$\sigma(f),\sigma(F)$分别为$f,F$的方差。
\end{theorem}
物理规律中许多现象都可以归结为这一数学事实，例如著名的\textbf{海森堡不确定性原理} (Heisenberg's uncertainty principle)\footnote{尽管广义的不确定性原理应该归结为算符或矩阵的不对易，但由于动量空间波函数是位置空间波函数的傅里叶变换，以上结果的确能够证明$\sigma_x \sigma_p\geq\frac{\hbar}{2}$}。读者可以自行验证，以上不等式取等于高斯函数。
\begin{proof}
    我们首先验证，$|f|^2$与$|F|^2$作为某个随机变量的概率密度函数是良定义的，从而其方差是良定义的。根据普朗歇尔定理（见\ref{sub:fourier_operation_properties}），\lr{
    \int_{-\infty}^{\infty}|f(t)|^2\,dt=\int_{-\infty}^{\infty}|\mathcal{F} f(\xi)|^2\,d\xi
}{
    \int_{-\infty}^{\infty}|f(t)|^2\,dt=\frac{1}{2\pi}\int_{-\infty}^{\infty}|\mathcal{F} f(\omega)|^2\,d\omega
}
    要使它们都是概率密度函数，我们就需要把$|f|^2,|F|^2$归一化，在频率形式的傅里叶变换下，这是自动成立的，而如果使用角频率来做傅里叶变换，则需要引入额外的归一化系数，因此下面仅使用频率。

    我们知道，如果随机变量$X$的概率密度函数为$f(x)$，并且其方差有定义（二阶中心矩存在），则方差为$\sigma^2=\int_{-\infty}^{\infty}(x-\mu)^2 f(x)\,dx$，其中$\mu=\int_{-\infty}^{\infty}x f(x)\,dx$为数学期望。
    
    我们假设$f$和$F$的数学期望为0，否则可以通过时移、频移定理将它调整为0，而不影响另一个域上信号的模值；又假设\[\int_{-\infty}^{\infty}t^2|f(t)|^2\,dt<\infty,\int_{-\infty}^{\infty}\xi^2|F(\xi)|^2\,d\xi<\infty\]
    即方差存在，则\begin{align*}
    4\pi^2 \sigma^2(f)\sigma^2(F)&=4\pi^2 \int_{-\infty}^{\infty}t^2|f(t)|^2\,dt\int_{-\infty}^{\infty}\xi^2|F(\xi)|^2\,d\xi\\
    &=\int_{-\infty}^{\infty}t^2|f(t)|^2\,dt\int_{-\infty}^{\infty}|2\pi \mathrm{i}\xi F(\xi)|^2\,d\xi\\
    &=\int_{-\infty}^{\infty}t^2|f(t)|^2\,dt\int_{-\infty}^{\infty}\left| \mathcal{F} [f'(t)](\xi)\right|^2\,d\xi&\text{(时域微分性质)}\\
    &=\int_{-\infty}^{\infty}t^2|f(t)|^2\,dt\int_{-\infty}^{\infty}| f'(t)|^2\,dt&\text{(普朗歇尔定理)}\\
    &\geq \left(\int_{-\infty}^{\infty}|tf(t)||f'(t)|\,dt\right)^2&\text{(柯西-施瓦兹不等式)}\\
    &=\frac{1}{4}\int_{-\infty}^{\infty}|f(t)|^2\,dt^2&\text{(分部积分)}\\
    &=\frac{1}{4}&\text{(归一化)}
    \end{align*}
    因此，$\sigma(f)\sigma(F)\geq \dfrac{1}{4\pi}$，得证。
\end{proof}

使用方差的定义可以轻松地得到一个有用的不等式，它给出了一种基于方差估计函数集中程度的方法，这就是\textbf{切比雪夫不等式} (Chebyshev's inequality)：
\begin{theorem}[切比雪夫不等式]
    设$f\in L^2(\mathbb{R} )$，则对于任意$a>0$，都有
    \[\int_{|t|\geq a}|f(t)|^2\,dt\leq \frac{1}{a^2}\int_{-\infty}^{\infty}t^2|f(t)|^2\,dt\]
    在概率论中它的表述是：设随机变量$X$的数学期望为$\mu$，方差为$\sigma^2$，则对于任意$k>0$，都有\[P(|X-\mu|\geq a)\leq \frac{\sigma^2}{a^2}\]
\end{theorem}
\begin{proof}
    由于$a>0$，有
    \begin{align*}
    \int_{|t|\geq a}|f(t)|^2\,dt&=\int_{|t|\geq a}\frac{t^2}{t^2}|f(t)|^2\,dt\\
    &\leq \frac{1}{a^2}\int_{|t|\geq a}t^2|f(t)|^2\,dt\\
    &\leq \frac{1}{a^2}\int_{-\infty}^{\infty}t^2|f(t)|^2\,dt
\end{align*}
\end{proof}

\begin{example}
    设$f$满足归一化条件且数学期望为0，取$a=2\sigma(f)$，则有\[\int_{|t|\geq 2\sigma(f)}|f(t)|^2\,dt\leq \frac{1}{4}\]
即$f$在区间$[-2\sigma(f),2\sigma(f)]$之外的能量不超过总能量的$1/4$。
\end{example}
可见，前面证明的$\sigma_f \sigma_F\geq 1/4\pi$又可以表述为时域上的“宽度”与频域上的“宽度”之积有下界，由此可以定义信号的时宽、带宽积：
\begin{definition}[时宽-带宽积]
    设$f\in L^1(\mathbb{R} )\cap L^2(\mathbb{R} )$，$F(\xi)=\int_{-\infty}^{\infty}f(t)\mathrm{e}^{-2\pi \mathrm{i}\xi t}\,dt$，且它们的方差存在，则$f$的\textbf{时宽}(time width)和\textbf{带宽}(bandwidth)定义分别为
    \[T=2\sigma(f),B=2\sigma(F)\]
    它们的乘积称为\textbf{时宽-带宽积} (time-bandwidth product)：
    \[TB(f)=4\sigma(f)\sigma(F)\]
\end{definition}

我们也可以根据时域、频域函数的半高宽或$1/20$高宽（1dB带宽）来估计它们的集中程度。

\subsection{*傅里叶变换的连续性}\label{sub:fourier_continuous}
施瓦兹函数类的另一个好处在于，它不仅保证了傅里叶变换下的封闭性，还使得傅里叶变换$\mathcal{F} :\mathcal{S} \to \mathcal{S} $是一个\textbf{线性拓扑自同构}，具体来说，这要求\begin{itemize}
    \item $\mathcal{F} $是线性的
    \item $\mathcal{F} $是双射（我们已经找到逆映射$\mathcal{F} ^{-1} $）
    \item $\mathcal{F} $及其逆映射$\mathcal{F} ^{-1} $均为连续映射。
\end{itemize}
我们已经证明过$\mathcal{F}$是线性的双射，下面来说明连续性。为此，我们需要给施瓦兹函数类赋予一个拓扑结构。尽管$\mathcal{S}\in L^1(\mathbb{R} )$，但$L^p$空间的范数并不能很好地刻画施瓦兹函数类的性质，因此我们需要引入一族\textbf{半范数} (seminorm)：
\begin{definition}[*施瓦兹函数类的半范数]
    对于$m,n\in\mathbb{N} $，定义施瓦兹函数类上的半范数
    \[\left\|\varphi\right\|_{m,n}=\sup_{x\in\mathbb{R} }|x|^m\left|\varphi^{(n)}(x)\right|\]
\end{definition}
在这个范数下$\mathcal{S} $构成\textbf{弗雷歇空间} (Fréchet space)，依范数收敛此时应该理解为：
\begin{definition}[*施瓦兹函数类中的收敛]
    设$\{\varphi_k\}\subset\mathcal{S} $，如果对于任意$m,n\in\mathbb{N} $，都有
    \[\lim_{k\to\infty}\left\|\varphi_k-\varphi\right\|_{m,n}=0\]
    则称$\varphi_k$在施瓦兹函数类中（依范数）收敛于$\varphi$，记为$\varphi_k\to \varphi$。
\end{definition}
有了这些定义，我们就可以说明傅里叶变换及其逆映射的连续性：
\begin{theorem}[*$\mathcal{S} $上傅里叶变换的连续性]
    设$\varphi_k\to \varphi$，则$\mathcal{F} \varphi_k\to \mathcal{F} \varphi$。
\end{theorem}
\begin{proof}
    不妨设$\varphi=\mathcal{F} \varphi=0$。由傅里叶变换的微分性质（见\ref{sub:general_distributions}），
    \begin{align*}
    \mathcal{F}\left[\left((-it)^n\varphi_k\right)^{(m)}\right](\omega)&=(\mathrm{i}\omega)^m \mathcal{F}\left[(-it)^n \varphi_k\right](\omega)\\
    &=(\mathrm{i}\omega)^m (\mathcal{F} \varphi_k)^{(n)}(\omega)\\
    \|\mathcal{F} \varphi_k\|_{m,n}&=\sup_{\omega\in\mathbb{R} }|\omega|^m\left|(\mathcal{F} \varphi_k)^{(n)}(\omega)\right|\\
    &=\sup_{\omega\in\mathbb{R} }\left|\mathcal{F}\left[((-it)^n \varphi_k)^{(m)}\right](\omega)\right|
\end{align*}
令$\phi_k=\left((-it)^n \varphi_k\right)^{(m)}\in\mathcal{S} $，则
\begin{align*}
    \|\mathcal{F} \varphi_k\|_{m,n}&=\sup_{\omega\in\mathbb{R} }\left|\mathcal{F} \phi_k(\omega)\right|\leq \int_{-\infty}^{\infty}|\phi_k(t)|\,dt
\end{align*}
问题转化为证明\[\int_{-\infty}^{\infty}|\phi_k(t)|\,dt\to 0\]
考虑利用$\varphi_k\in\mathcal{S} ,\|\varphi_k\|_{m,n}\to 0,\forall m,n\in\mathbb{N}$给出$|\phi_k(t)|$的估计。利用莱布尼兹求导法则，有：
\begin{align*}
    |\phi_k(t)|&=\left|\left((-it)^n \varphi_k\right)^{(m)}(t)\right|\\
&=\left|\sum_{j=0}^{m}\binom{m}{j}(t^n)^{(j)}\varphi_k^{(m-j)}(t)\right|\\
    &\leq \sum_{j=0}^{\min(m,n)}\binom{m}{j}\frac{n!}{(n-j)!}|t|^{n-j}\left|\varphi_k^{(m-j)}(t)\right|
\end{align*}
由于$\|\varphi_k\|_{n-j+2,m-j}\to 0$，有：$\forall \epsilon>0,\exists N\in\mathbb{N},k>N\implies \|\varphi_k\|_{n-j+2,m-j}<\epsilon$，于是
\begin{align*}
    |\phi_k(t)|&\leq \sum_{j=0}^{\min(m,n)}\binom{m}{j}\frac{n!}{(n-j)!}\frac{\|\varphi_k\|_{n-j+2,m-j}}{|t|^2}\\
    &<\frac{C\epsilon}{|t|^2},C=\sum_{j=0}^{\min(m,n)}\binom{m}{j}\frac{n!}{(n-j)!}
\end{align*}
因此有\begin{align*}
    \int_{|t|\geq 1}|\phi_k(t)|\,dt&\leq \int_{|t|\geq 1}\frac{C\epsilon}{|t|^2}\,dt\\
    &=2C\epsilon\int_{1}^{\infty}\frac{1}{t^2}\,dt=2C\epsilon
\end{align*}
而在$|t|<1$时，\begin{align*}
    \sum_{j=0}^{\min(m,n)}\binom{m}{j}\frac{n!}{(n-j)!}|t|^{n-j}\left|\varphi_k^{(m-j)}(t)\right|&\leq \sum_{j=0}^{\min(m,n)}\binom{m}{j}\frac{n!}{(n-j)!}\|\varphi_k\|_{0,m-j}\\
    &\leq C'\max_{0\leq j\leq \min(m,n)}\|\varphi_k\|_{0,m-j}<\epsilon
\end{align*}
即$\phi_k$一致收敛于0，因此有
\[\int_{|t|<1}|\phi_k(t)|\,dt\to 0\]
综上，\[0\leq\|\mathcal{F} \varphi_k\|_{m,n}\leq\int_{-\infty}^{\infty}|\phi_k(t)|\,dt\to 0\]
根据夹逼定理，$\|\mathcal{F} \varphi_k\|_{m,n}\to 0$，得证。傅里叶逆变换同理。
\end{proof}

其实，如果我们考虑$L^1(\mathbb{R})\cap L^2(\mathbb{R})$中傅里叶变换的连续性，可以得到类似的结论，但证明要简单得多，因为$L^1(\mathbb{R})\cap L^2(\mathbb{R})$中有普朗歇尔定理（见\ref{sub:fourier_operation_properties}）。
\begin{theorem}[$L^1(\mathbb{R})\cap L^2(\mathbb{R})$上傅里叶变换的连续性]
    设$f_k$在$L^2$范数下收敛到$f$，则$\mathcal{F} f_k$在$L^2$范数下收敛到$\mathcal{F} f$。
\end{theorem}
\begin{remark}
    尽管$L^1(\mathbb{R})\cap L^2(\mathbb{R})$是比$\mathcal{S} $更大的空间，但这个定理也有一定的局限性：其一，依$L^2$范数收敛不能保证依半范数收敛，即证得的收敛性会减弱；其二，$L^1(\mathbb{R})\cap L^2(\mathbb{R})$并不是傅里叶变换下的封闭空间，因此我们无法保证傅里叶逆变换的连续性。
\end{remark}
\begin{proof}
    由普朗歇尔定理，有
    \lr{
    \left\|\mathcal{F} f_k-\mathcal{F} f\right\|_2^2&=\int_{-\infty}^{\infty}\left|\mathcal{F} f_k(\xi)-\mathcal{F} f(\xi)\right|^2\,d\xi\\
    &=\int_{-\infty}^{\infty}\left|f_k(t)-f(t)\right|^2\,dt\\
    &=\left\|f_k-f\right\|_2^2
    }{
    \left\|\mathcal{F} f_k-\mathcal{F} f\right\|_2^2&=\int_{-\infty}^{\infty}\left|\mathcal{F} f_k(\omega)-\mathcal{F} f(\omega)\right|^2\,d\omega\\
    &=2\pi\int_{-\infty}^{\infty}\left|f_k(t)-f(t)\right|^2\,dt\\
    &=2\pi\left\|f_k-f\right\|_2^2
    }
    因此，$\left\|f_k-f\right\|_2\to 0\Rightarrow \left\|\mathcal{F} f_k-\mathcal{F} f\right\|_2\to 0$，得证。
\end{proof}

另外，在弱收敛的意义下，傅里叶变换同样是$\mathcal{T} =\mathcal{S} ' $上的线性拓扑自同构。在\ref{sub:operation_of_distribution}中，我们已经证明过分布的傅里叶变换的连续性：\begin{theorem}[$\mathcal{T}$上傅里叶变换的连续性]
    设$T_k\to T$，则$\mathcal{F} T_k\to \mathcal{F} T$。
\end{theorem}
